\chapter{Resumen}
Kubernetes no impone un plugin de red específico, lo que obliga a los
equipos a seleccionar un \emph{Container Network Interface} (CNI) sin
métricas comparativas objetivas. Esta situación conduce a
sobredimensionamiento de infraestructura (CPU/RAM) y a posturas de
seguridad incompletas. El presente trabajo construye evidencia empírica
reproducible —infraestructura como código, benchmarks automatizados y
modelo de decisión multicriterio (MCDA)— para guiar esa selección.
Se evalúan Flannel, Calico, Cilium y Antrea en un clúster K3s en alta
disponibilidad sobre DigitalOcean, midiendo \emph{throughput} TCP,
latencia de establecimiento, retransmisiones y overhead de CPU/RAM del
agente CNI, así como el impacto de políticas de red Zero Trust.
Los resultados alimentan un modelo MCDA parametrizable por perfil de
uso (Fintech/seguridad, Streaming/rendimiento e IoT/recursos limitados) y un prototipo
web que expone las recomendaciones.
Los hallazgos principales muestran que ningún plugin resultó superior
en todos los criterios: Calico alcanzó el mayor throughput TCP
(1804.40 Mbps) y la menor latencia de establecimiento, mientras Flannel
registró la menor tasa de retransmisión y huella de recursos. Se
comprobó además que Flannel ---CNI predeterminado de K3s--- acepta
NetworkPolicies en la API de Kubernetes pero no las aplica en su plano
de datos, configurando un falso positivo de seguridad, y que Cilium fue
el único plugin capaz de aplicar políticas de capa 7 (HTTP). En
coherencia con esa ausencia de dominancia, el modelo arroja una
recomendación distinta por perfil: Cilium para Fintech, Calico para
Streaming y Flannel para IoT, este último por un margen estrecho.

\vspace{1cm}
\textbf{Palabras clave:} Kubernetes, CNI, NetworkPolicy, Zero Trust, MCDA, QoS, eBPF.